% Options for packages loaded elsewhere
\PassOptionsToPackage{unicode}{hyperref}
\PassOptionsToPackage{hyphens}{url}
\PassOptionsToPackage{dvipsnames,svgnames,x11names}{xcolor}
%
\documentclass[
  a4paper,
  11pt]{article}
\usepackage{amsmath,amssymb}
\usepackage{iftex}
\ifPDFTeX
  \usepackage[T1]{fontenc}
  \usepackage[utf8]{inputenc}
  \usepackage{textcomp} % provide euro and other symbols
\else % if luatex or xetex
  \usepackage{unicode-math} % this also loads fontspec
  \defaultfontfeatures{Scale=MatchLowercase}
  \defaultfontfeatures[\rmfamily]{Ligatures=TeX,Scale=1}
\fi
\usepackage{lmodern}
\ifPDFTeX\else
  % xetex/luatex font selection
\fi
% Use upquote if available, for straight quotes in verbatim environments
\IfFileExists{upquote.sty}{\usepackage{upquote}}{}
\IfFileExists{microtype.sty}{% use microtype if available
  \usepackage[]{microtype}
  \UseMicrotypeSet[protrusion]{basicmath} % disable protrusion for tt fonts
}{}
\makeatletter
\@ifundefined{KOMAClassName}{% if non-KOMA class
  \IfFileExists{parskip.sty}{%
    \usepackage{parskip}
  }{% else
    \setlength{\parindent}{0pt}
    \setlength{\parskip}{6pt plus 2pt minus 1pt}}
}{% if KOMA class
  \KOMAoptions{parskip=half}}
\makeatother
\usepackage{xcolor}
\usepackage[margin=2.4cm]{geometry}
\usepackage{pdflscape}
\usepackage{longtable,booktabs,array}
\usepackage{calc} % for calculating minipage widths
% Correct order of tables after \paragraph or \subparagraph
\usepackage{etoolbox}
\makeatletter
\patchcmd\longtable{\par}{\if@noskipsec\mbox{}\fi\par}{}{}
\makeatother
% Allow footnotes in longtable head/foot
\IfFileExists{footnotehyper.sty}{\usepackage{footnotehyper}}{\usepackage{footnote}}
\makesavenoteenv{longtable}
\setlength{\emergencystretch}{3em} % prevent overfull lines
\providecommand{\tightlist}{%
  \setlength{\itemsep}{0pt}\setlength{\parskip}{0pt}}
\setcounter{secnumdepth}{5}
\ifLuaTeX
\usepackage[bidi=basic]{babel}
\else
\usepackage[bidi=default]{babel}
\fi
\babelprovide[main,import]{british}
% get rid of language-specific shorthands (see #6817):
\let\LanguageShortHands\languageshorthands
\def\languageshorthands#1{}
\usepackage{xurl}
\usepackage{enumitem}
\usepackage{ragged2e}
\setlist{nosep,topsep=0.45em}
\setlength{\emergencystretch}{3em}
\urlstyle{same}
\ifLuaTeX
  \usepackage{selnolig}  % disable illegal ligatures
\fi
\IfFileExists{bookmark.sty}{\usepackage{bookmark}}{\usepackage{hyperref}}
\IfFileExists{xurl.sty}{\usepackage{xurl}}{} % add URL line breaks if available
\urlstyle{same}
\hypersetup{
  pdftitle={Artificial Intelligence for Earth Observation},
  pdfauthor={Stefano Infusini},
  pdflang={en-GB},
  colorlinks=true,
  linkcolor={black},
  filecolor={Maroon},
  citecolor={blue},
  urlcolor={blue},
  pdfcreator={LaTeX via pandoc}}

\title{Artificial Intelligence for Earth Observation}
\usepackage{etoolbox}
\makeatletter
\providecommand{\subtitle}[1]{% add subtitle to \maketitle
  \apptocmd{\@title}{\par {\large #1 \par}}{}{}
}
\makeatother
\subtitle{Current research landscape, scientific development, open
problems, and MSc thesis directions}
\author{Stefano Infusini}
\date{August 2026}

\begin{document}
\pagenumbering{gobble}
\maketitle

\thispagestyle{empty}
\vfill
\begin{center}
\begin{minipage}{0.86\textwidth}
\centering
\textbf{Purpose:} Technical orientation for discussion of feasible MSc thesis directions\\[0.6em]
\textbf{Literature cut-off:} 9 August 2026
\end{minipage}
\end{center}
\vfill
\clearpage
\pagenumbering{roman}

\renewcommand*\contentsname{Contents}
{
\hypersetup{linkcolor=black}
\setcounter{tocdepth}{3}
\tableofcontents
}
\clearpage
\pagenumbering{arabic}
\textbf{Evidence note:} The overview prioritizes original papers,
official benchmark reports, and ESA/NASA sources. Because the field is
moving quickly, it includes both peer-reviewed work and clearly
identifiable recent preprints; claims of state of the art from
individual model papers are treated as provisional unless supported by
independent benchmarks.

\begin{center}\rule{0.5\linewidth}{0.5pt}\end{center}

\hypertarget{executive-synthesis}{%
\section*{Executive synthesis}\label{executive-synthesis}}
\addcontentsline{toc}{section}{Executive synthesis}

Artificial Intelligence for Earth Observation (AI4EO) is not a single
task or model family. It is the research area concerned with converting
measurements of the Earth---optical, multispectral, hyperspectral,
thermal, radar, LiDAR, elevation, meteorological, in-situ, and textual
data---into reliable information about the state and evolution of the
planet.

The field has developed through five broad transitions:

\begin{enumerate}
\def\labelenumi{\arabic{enumi}.}
\tightlist
\item
  \textbf{From physically designed features and statistical classifiers
  to learned representations.} Early operational pipelines depended on
  spectral indices, texture descriptors, dimensionality reduction,
  maximum-likelihood classifiers, Support Vector Machines, Random
  Forests, and expert-designed rules.
\item
  \textbf{From pixel-wise models to spatial and spectral--spatial deep
  learning.} CNNs, fully convolutional networks, and encoder--decoder
  models made it possible to learn spatial context, perform dense
  mapping, and process very-high-resolution imagery.
\item
  \textbf{From single images to multimodal and temporal learning.}
  Research began to exploit optical--SAR complementarity, satellite
  image time series, phenology, change, irregular observations, and
  auxiliary geospatial variables.
\item
  \textbf{From task-specific supervised models to self-supervised
  pretraining and EO foundation models.} Large unlabelled archives are
  now used to pretrain encoders that can be adapted to multiple tasks
  with fewer labels.
\item
  \textbf{From benchmark accuracy to capability, reliability, and
  deployment.} The most important current questions concern geographic
  and sensor transfer, missing modalities, trustworthy uncertainty,
  physical validity, reproducible evaluation, computational efficiency,
  natural-language interaction, and tool-using agents.
\end{enumerate}

The main conclusion of the 2024--2026 literature is deliberately more
cautious than the current foundation-model enthusiasm:

\begin{quote}
\textbf{There is no single state-of-the-art AI4EO model.} Performance
depends on the sensor, spectral bands, spatial and temporal resolution,
task, label budget, adaptation protocol, and target geography.
\end{quote}

This conclusion is supported by several independent evaluations.
\href{https://arxiv.org/abs/2412.04204}{PANGAEA} found that geospatial
foundation models do not consistently outperform supervised U-Net and
ViT baselines. \href{https://arxiv.org/html/2511.15658v1}{GEO-Bench-2}
found that EO-specific pretraining is particularly valuable when
multispectral or temporal information matters, while large natural-image
models remain highly competitive on high-resolution RGB tasks. A 2026
audit of 152 foundation-model papers found major inconsistencies in
evaluation and concluded that the published record does not presently
support a universal ranking of models
(\href{https://arxiv.org/html/2605.12678v1}{Corley et al., 2026}).

For an MSc thesis, this is good news. A meaningful contribution does
\textbf{not} require pretraining another giant model. A stronger and
more feasible thesis can use open pretrained models to investigate a
precise EO problem such as:

\begin{itemize}
\tightlist
\item
  transfer to unseen cities, regions, seasons, or sensors;
\item
  robustness when one sensor is cloudy, delayed, or unavailable;
\item
  uncertainty calibration under distribution shift;
\item
  physically consistent reconstruction or downscaling;
\item
  fair, leakage-resistant evaluation of EO foundation models;
\item
  reliability of language agents that select and process EO products.
\end{itemize}

Given Stefano's background in AI/CV, signal and image processing, NLP,
HPC, and earlier Landsat--Sentinel-2 land-surface-temperature work, the
strongest provisional direction is:

\begin{quote}
\textbf{Uncertainty-aware and physically consistent LST downscaling
using EO foundation-model representations, evaluated across cities,
seasons, and sensors.}
\end{quote}

This is not a final topic selection. It is the most promising
intersection identified so far between prior experience, current
research gaps, open data, tractable compute, and a defensible
experimental protocol.

\begin{center}\rule{0.5\linewidth}{0.5pt}\end{center}

\hypertarget{how-to-organize-the-ai4eo-field}{%
\section{How to organize the AI4EO
field}\label{how-to-organize-the-ai4eo-field}}

Many reviews mix together data types, computer-vision tasks, scientific
applications, and methodological research trends. They are different
axes.

\begin{longtable}[]{@{}
  >{\raggedright\arraybackslash}p{(\columnwidth - 4\tabcolsep) * \real{0.3333}}
  >{\raggedright\arraybackslash}p{(\columnwidth - 4\tabcolsep) * \real{0.3333}}
  >{\raggedright\arraybackslash}p{(\columnwidth - 4\tabcolsep) * \real{0.3333}}@{}}
\toprule\noalign{}
\begin{minipage}[b]{\linewidth}\raggedright
Axis
\end{minipage} & \begin{minipage}[b]{\linewidth}\raggedright
Representative categories
\end{minipage} & \begin{minipage}[b]{\linewidth}\raggedright
What it answers
\end{minipage} \\
\midrule\noalign{}
\endhead
\bottomrule\noalign{}
\endlastfoot
\textbf{Observation modality} & RGB/optical, multispectral,
hyperspectral, thermal, SAR, LiDAR/DEM, meteorology, in-situ data, text
& What is measured, by which physical process? \\
\textbf{AI task} & Classification, semantic/instance/panoptic
segmentation, object detection, regression/retrieval, change detection,
forecasting, reconstruction, super-resolution, retrieval, VQA & What
mathematical output must the system produce? \\
\textbf{Scientific/application domain} & Agriculture, forestry,
biodiversity, disasters, climate, urban systems, ocean, cryosphere,
atmosphere, water, geology & What Earth-system or societal question is
being addressed? \\
\textbf{Research branch} & Foundation models, multimodal fusion,
spatio-temporal learning, generalization, trustworthy/physics-aware AI,
vision-language and agents & What reusable scientific or methodological
problem is being studied? \\
\end{longtable}

A thesis topic should normally select one element from each axis. For
example:

\begin{quote}
\textbf{Sentinel-1 + Sentinel-2} (modalities) + \textbf{flood
segmentation} (task) + \textbf{disaster response} (application) +
\textbf{missing-modality robustness and uncertainty under event shift}
(research problem).
\end{quote}

This formulation is much sharper than a topic such as ``use Transformers
for satellite images,'' because it specifies the scientific question and
the conditions under which the contribution can be evaluated.

\hypertarget{task-taxonomy}{%
\subsection{Task taxonomy}\label{task-taxonomy}}

The main AI tasks in EO are:

\begin{itemize}
\tightlist
\item
  \textbf{Scene or patch classification:} assign one or multiple labels
  to an image patch.
\item
  \textbf{Pixel-wise classification / semantic segmentation:} assign a
  class to every pixel; land-cover mapping is usually formulated this
  way when a spatial map is required.
\item
  \textbf{Object detection:} localize countable entities with boxes or
  oriented boxes, such as ships, aircraft, vehicles, buildings, or wind
  turbines.
\item
  \textbf{Instance or panoptic segmentation:} separate individual
  objects or parcels while also assigning semantic classes.
\item
  \textbf{Regression and geophysical retrieval:} estimate continuous
  quantities such as LST, biomass, soil moisture, crop yield,
  chlorophyll, or atmospheric variables.
\item
  \textbf{Change detection:} identify where, what, and sometimes when a
  meaningful surface change occurred.
\item
  \textbf{Time-series modelling and forecasting:} infer phenology,
  monitor event evolution, fill gaps, or predict future states.
\item
  \textbf{Reconstruction and enhancement:} cloud removal, denoising, gap
  filling, super-resolution, data fusion, and cross-modal generation.
\item
  \textbf{Retrieval, captioning, visual question answering, and
  grounding:} connect imagery with language and external knowledge.
\end{itemize}

Land-cover classification and object detection are therefore different
problem formulations. In land-cover mapping, the output is generally a
dense categorical field. In object detection, the output is a set of
localized object instances. Neither task is itself a complete ``research
branch'': both can be studied using multimodal learning, foundation
models, domain adaptation, uncertainty estimation, or other branches.

\begin{center}\rule{0.5\linewidth}{0.5pt}\end{center}

\hypertarget{scientific-development-of-ai-for-eo}{%
\section{Scientific development of AI for
EO}\label{scientific-development-of-ai-for-eo}}

\hypertarget{developmental-timeline}{%
\subsection{Developmental timeline}\label{developmental-timeline}}

\begin{longtable}[]{@{}
  >{\raggedright\arraybackslash}p{(\columnwidth - 6\tabcolsep) * \real{0.2500}}
  >{\raggedright\arraybackslash}p{(\columnwidth - 6\tabcolsep) * \real{0.2500}}
  >{\raggedright\arraybackslash}p{(\columnwidth - 6\tabcolsep) * \real{0.2500}}
  >{\raggedright\arraybackslash}p{(\columnwidth - 6\tabcolsep) * \real{0.2500}}@{}}
\toprule\noalign{}
\begin{minipage}[b]{\linewidth}\raggedright
Period
\end{minipage} & \begin{minipage}[b]{\linewidth}\raggedright
Dominant scientific idea
\end{minipage} & \begin{minipage}[b]{\linewidth}\raggedright
What changed in EO
\end{minipage} & \begin{minipage}[b]{\linewidth}\raggedright
Principal limitation left behind
\end{minipage} \\
\midrule\noalign{}
\endhead
\bottomrule\noalign{}
\endlastfoot
\textbf{Before \textasciitilde2005} & Physical models, statistical
pattern recognition, expert-designed features & Spectral indices,
transforms, per-pixel classifiers, rule-based retrieval, and
signal-processing pipelines encoded domain knowledge directly & Limited
ability to learn complex spatial patterns; extensive feature
engineering \\
\textbf{\textasciitilde2005--2014} & Classical machine learning & SVMs,
Random Forests, boosting, object-based image analysis, kernel methods,
and sparse representations improved nonlinear classification with
comparatively small labelled datasets & Representations remained largely
hand-designed; spatial context was often shallow \\
\textbf{\textasciitilde2014--2018} & Deep representation learning &
Autoencoders, DBNs, 1-D/2-D/3-D CNNs, transfer learning, FCNs, and
encoder--decoder models learned spectral, spatial, and joint
spectral--spatial features & Label hunger, patch-level evaluation,
limited geographic transfer, and dependence on ImageNet/RGB
assumptions \\
\textbf{\textasciitilde2018--2021} & Dataset and task scaling &
BigEarthNet, SEN12MS, SpaceNet, xView, So2Sat, PASTIS, and other
benchmarks supported dense mapping, multimodal fusion, and time-series
learning & Datasets remained geographically and semantically fragmented;
random splits often overstated transfer \\
\textbf{\textasciitilde2020--2023} & Self-supervision and Transformers &
Contrastive learning, masked modelling, temporal attention, and
sensor-aware pretraining exploited massive unlabelled archives & Models
often remained tied to fixed sensors, band configurations, resolutions,
or tasks \\
\textbf{\textasciitilde2023--2025} & EO foundation models & Prithvi,
CROMA, DOFA, SkySense, AnySat, and related models pursued reusable,
multimodal, multi-temporal representations & Comparisons were
inconsistent; universal superiority was not established \\
\textbf{\textasciitilde2025--2026} & Capability, generation, reasoning,
and operations & TerraMind, OlmoEarth, RAMEN, vision-language systems,
agents, global embeddings, lightweight variants, and on-orbit
demonstrations expanded the target from feature extraction to flexible
Earth intelligence & Reliability, physical faithfulness, cross-domain
transfer, evaluation, provenance, and deployment remain open \\
\end{longtable}

The classical-to-deep-learning transition is documented in the 2017
review by \href{https://arxiv.org/abs/1710.03959}{Zhu et al.}. That
paper is useful historically because it captures the moment at which
remote sensing was moving from hand-crafted features to learned
spectral--spatial representations, while already warning against
treating deep learning as a domain-free black box.

\hypertarget{the-dataset-transition}{%
\subsection{The dataset transition}\label{the-dataset-transition}}

Deep learning became scientifically useful in EO only when the community
began constructing larger, georeferenced datasets and benchmarks.

\begin{itemize}
\tightlist
\item
  \href{https://arxiv.org/abs/1902.06148}{BigEarthNet} introduced
  590,326 multi-label Sentinel-2 patches in its original release,
  demonstrating the importance of domain-specific large-scale
  pretraining rather than relying only on ImageNet.
\item
  \href{https://arxiv.org/abs/1906.07789}{SEN12MS} provided 180,662
  globally distributed Sentinel-1/Sentinel-2/MODIS triplets, helping
  establish radar--optical learning and multimodal land-cover research.
\item
  \href{https://arxiv.org/abs/2107.07933}{PASTIS} made satellite image
  time series and agricultural parcel segmentation a central benchmark
  problem.
\item
  \href{https://arxiv.org/abs/2211.07044}{SSL4EO-S12} explicitly
  constructed a global, multi-seasonal Sentinel-1/2 corpus for
  self-supervised learning.
\item
  \href{https://arxiv.org/html/2402.12095v2}{Major TOM} moved toward
  expandable, AI-ready global EO data and, later, openly distributed
  global embedding products.
\end{itemize}

This history matters because model progress cannot be separated from
data design. Label ontologies, spatial coverage, temporal sampling,
cloud filtering, co-registration, and train/test geography all determine
what a model can learn and what a benchmark score means.

\hypertarget{the-self-supervised-and-transformer-transition}{%
\subsection{The self-supervised and Transformer
transition}\label{the-self-supervised-and-transformer-transition}}

EO is unusually well suited to self-supervised learning: unlabelled
measurements are abundant, while high-quality labels require expensive
expert interpretation and may exist only for particular places or dates.

\href{https://openaccess.thecvf.com/content/ICCV2021/papers/Manas_Seasonal_Contrast_Unsupervised_Pre-Training_From_Uncurated_Remote_Sensing_Data_ICCV_2021_paper.pdf}{Seasonal
Contrast (SeCo)} used repeated observations of the same locations to
learn seasonally robust representations.
\href{https://proceedings.neurips.cc/paper_files/paper/2022/hash/01c561df365429f33fcd7a7faa44c985-Abstract-Conference.html}{SatMAE}
adapted masked autoencoding to temporal and multispectral satellite
imagery through temporal and spectral positional information. These
works changed the scientific question from ``How do we label enough
pixels?'' to ``Which invariant and predictive structures can be learned
directly from the observation archive?''

\hypertarget{the-foundation-model-transition}{%
\subsection{The foundation-model
transition}\label{the-foundation-model-transition}}

The term \emph{EO foundation model} is now used for a pretrained model
intended to be adapted across several EO tasks, regions, sensors, or
label regimes. Representative developments include:

\begin{longtable}[]{@{}
  >{\raggedright\arraybackslash}p{(\columnwidth - 6\tabcolsep) * \real{0.2308}}
  >{\raggedleft\arraybackslash}p{(\columnwidth - 6\tabcolsep) * \real{0.3077}}
  >{\raggedright\arraybackslash}p{(\columnwidth - 6\tabcolsep) * \real{0.2308}}
  >{\raggedright\arraybackslash}p{(\columnwidth - 6\tabcolsep) * \real{0.2308}}@{}}
\toprule\noalign{}
\begin{minipage}[b]{\linewidth}\raggedright
Model
\end{minipage} & \begin{minipage}[b]{\linewidth}\raggedleft
Year
\end{minipage} & \begin{minipage}[b]{\linewidth}\raggedright
Main scientific idea
\end{minipage} & \begin{minipage}[b]{\linewidth}\raggedright
Important limitation or question
\end{minipage} \\
\midrule\noalign{}
\endhead
\bottomrule\noalign{}
\endlastfoot
\href{https://openaccess.thecvf.com/content/ICCV2021/papers/Manas_Seasonal_Contrast_Unsupervised_Pre-Training_From_Uncurated_Remote_Sensing_Data_ICCV_2021_paper.pdf}{SeCo}
& 2021 & Temporal/seasonal contrastive pretraining & Primarily optical;
not a universal sensor model \\
\href{https://proceedings.neurips.cc/paper_files/paper/2022/hash/01c561df365429f33fcd7a7faa44c985-Abstract-Conference.html}{SatMAE}
& 2022 & Masked modelling for temporal and multispectral imagery & Fixed
input assumptions and adaptation choices still matter \\
\href{https://arxiv.org/abs/2304.14065}{Presto} & 2023 & Lightweight
pretraining for multimodal pixel time series & Does not model fine
spatial structure in the same way as image encoders \\
\href{https://arxiv.org/abs/2311.00566}{CROMA} & 2023 & Contrastive
radar--optical masked autoencoding & Focused on Sentinel-1/2
radar--optical representations \\
\href{https://arxiv.org/abs/2412.02732}{Prithvi-EO-2.0} & 2024 & Global
HLS multi-temporal masked pretraining with temporal/location embeddings
& 300M/600M variants remain expensive relative to small task-specific
models \\
\href{https://arxiv.org/abs/2403.15356}{DOFA} & 2024 &
Wavelength-conditioned dynamic model for heterogeneous sensors &
Wavelength metadata does not solve all non-spectral modality and
resolution shifts \\
\href{https://arxiv.org/abs/2312.10115}{SkySense} & 2024 & Billion-scale
multimodal, temporal, and geo-context pretraining & Size,
reproducibility, and adaptation complexity \\
\href{https://arxiv.org/abs/2412.14123}{AnySat} & 2025 & One JEPA-based
model across heterogeneous resolutions, scales, modalities, and datasets
& A single shared model must still negotiate negative transfer and
task-specific needs \\
\href{https://arxiv.org/html/2504.11171v5}{TerraMind} & 2025 &
Any-to-any generative multimodal modelling and ``thinking in
modalities'' & Generated modalities must be tested for physical
faithfulness and downstream bias \\
\href{https://arxiv.org/abs/2511.13655}{OlmoEarth} & 2026 & Stable
latent modelling for multimodal, spatio-temporal EO & Claims across
heterogeneous benchmarks still depend on evaluation protocol \\
\href{https://arxiv.org/abs/2512.05025}{RAMEN} & 2026 & Sensor-agnostic,
resolution-adjustable representation and explicit detail/compute
trade-off & New-model results require broader independent
reproduction \\
\end{longtable}

The trajectory is clear: single-sensor encoders are becoming
multisensor, temporal, resolution-aware, and sometimes generative. The
unresolved question is not whether such models are promising. It is
\textbf{when their learned representation transfers more reliably or
efficiently than a well-tuned smaller baseline}.

\hypertarget{the-evaluation-transition}{%
\subsection{The evaluation transition}\label{the-evaluation-transition}}

The scientific maturation of the field is visible in its benchmarks:

\begin{itemize}
\tightlist
\item
  \href{https://proceedings.neurips.cc/paper_files/paper/2023/hash/a0644215d9cff6646fa334dfa5d29c5a-Abstract-Datasets_and_Benchmarks.html}{GEO-Bench}
  standardized six classification and six segmentation tasks.
\item
  \href{https://arxiv.org/abs/2401.04464}{PhilEO Bench} evaluated label
  efficiency using a global Sentinel-2 testbed and common decoder.
\item
  \href{https://arxiv.org/abs/2412.04204}{PANGAEA} expanded evaluation
  across tasks, resolutions, modalities, temporalities, and geographies.
\item
  \href{https://arxiv.org/html/2511.15658v1}{GEO-Bench-2} shifted from a
  single aggregate leaderboard toward capability-oriented evaluation.
\item
  \href{https://arxiv.org/html/2605.12678v1}{No One Knows the State of
  the Art in Geospatial Foundation Models} audited the literature itself
  and exposed a reproducibility and comparability crisis.
\end{itemize}

This is a transition from model-centric research---``our encoder gains
1.2 points''---to scientific evaluation---``which capability improved,
under what distribution shift, with which labels, preprocessing,
decoder, and compute?''

\begin{center}\rule{0.5\linewidth}{0.5pt}\end{center}

\hypertarget{current-core-research-branches}{%
\section{Current core research
branches}\label{current-core-research-branches}}

The following six branches provide a stable taxonomy for the main
scientific directions. Efficient deployment and data/benchmark
infrastructure are treated afterward as cross-cutting enabling layers.

\hypertarget{large-scale-pretraining-and-eo-foundation-models}{%
\subsection{Large-scale pretraining and EO foundation
models}\label{large-scale-pretraining-and-eo-foundation-models}}

\textbf{Central objective:} Learn reusable representations from large
unlabelled EO archives so that downstream models require fewer labels
and less task-specific training.

Current interests include:

\begin{itemize}
\tightlist
\item
  contrastive, masked, predictive, generative, and latent-space
  objectives;
\item
  sensor-specific versus sensor-agnostic encoders;
\item
  spatial, spectral, temporal, and geographic positional information;
\item
  frozen encoders, linear probing, full fine-tuning, adapters, LoRA, and
  prompt tuning;
\item
  scaling laws for model size, data volume, resolution, and modality
  diversity;
\item
  pretrained embeddings as a reusable data product;
\item
  model selection: predicting which encoder and adaptation strategy suit
  a target task.
\end{itemize}

\textbf{Open scientific problem:} Pretraining at scale does not
guarantee useful invariances. A model may learn location, season, sensor
artefacts, or land-cover prevalence rather than transferable
Earth-process structure.

\textbf{State of evidence:}

\begin{itemize}
\tightlist
\item
  PANGAEA found that current GFMs do not consistently beat supervised
  models.
\item
  GEO-Bench-2 found stronger benefits on multispectral tasks than on
  RGB-only tasks.
\item
  The 2026 audit found that inconsistent protocols make universal
  rankings scientifically unsupported.
\item
  A 2026 cross-region agriculture study found that Prithvi, SpectralGPT,
  and SatMAE all degraded sharply on unseen US states and tended to miss
  rare crops (\href{https://arxiv.org/abs/2606.29664}{Shang et al.,
  2026}).
\end{itemize}

\textbf{High-value research questions:}

\begin{enumerate}
\def\labelenumi{\arabic{enumi}.}
\tightlist
\item
  What part of downstream performance comes from pretraining data,
  architecture, model scale, or adaptation?
\item
  Which EO-specific capabilities---spectral, temporal, radar,
  geographic---are actually present in the representation?
\item
  When do frozen features suffice, and when is full fine-tuning
  necessary?
\item
  Can small models preserve the useful capabilities of large models?
\item
  Can model selection be performed without fine-tuning every candidate?
\end{enumerate}

\hypertarget{multimodal-and-multisensor-learning}{%
\subsection{Multimodal and multisensor
learning}\label{multimodal-and-multisensor-learning}}

\textbf{Central objective:} Combine complementary measurements of the
same Earth system.

Typical combinations include:

\begin{itemize}
\tightlist
\item
  optical + SAR;
\item
  multispectral + hyperspectral;
\item
  optical/SAR + DEM or LiDAR;
\item
  satellite imagery + weather/reanalysis;
\item
  EO imagery + in-situ measurements;
\item
  imagery + maps, coordinates, land-cover products, or text.
\end{itemize}

The research problem is more difficult than concatenating channels.
Different sensors measure different physical quantities, have different
noise models, resolutions, acquisition geometries, revisit schedules,
and missingness mechanisms.

Current interests include:

\begin{itemize}
\tightlist
\item
  early, intermediate, late, and token-level fusion;
\item
  shared versus modality-specific encoders;
\item
  cross-modal contrastive learning and masked reconstruction;
\item
  asynchronous fusion of non-simultaneous observations;
\item
  missing-modality training and inference;
\item
  cross-modal generation and imputation;
\item
  fusion without destructive resampling to a single nominal resolution.
\end{itemize}

Representative models include
\href{https://arxiv.org/abs/2311.00566}{CROMA},
\href{https://arxiv.org/abs/2412.14123}{AnySat},
\href{https://arxiv.org/html/2504.11171v5}{TerraMind},
\href{https://arxiv.org/abs/2511.13655}{OlmoEarth}, and
\href{https://arxiv.org/abs/2512.05025}{RAMEN}.

\textbf{Open scientific problem:} A multimodal model trained on
perfectly aligned data may fail in deployment precisely when
multimodality is most valuable---for example, when optical imagery is
cloudy, acquisitions are asynchronous, or one sensor is unavailable.

\hypertarget{spatio-temporal-modelling-and-change-analysis}{%
\subsection{Spatio-temporal modelling and change
analysis}\label{spatio-temporal-modelling-and-change-analysis}}

\textbf{Central objective:} Model the Earth as a dynamic system rather
than a collection of independent image chips.

Research topics include:

\begin{itemize}
\tightlist
\item
  bi-temporal and semantic change detection;
\item
  multi-temporal land-cover and crop mapping;
\item
  phenology and seasonal dynamics;
\item
  event onset, evolution, and recovery;
\item
  irregular sampling and cloud-related gaps;
\item
  long-context satellite image time series;
\item
  surface, weather, climate, and hazard forecasting;
\item
  continuous spatio-temporal representations.
\end{itemize}

\href{https://arxiv.org/abs/2107.07933}{PASTIS} showed the value of
temporal attention for agricultural parcel mapping.
\href{https://openaccess.thecvf.com/content/CVPR2021W/EarthVision/papers/Requena-Mesa_EarthNet2021_A_Large-Scale_Dataset_and_Challenge_for_Earth_Surface_Forecasting_CVPRW_2021_paper.pdf}{EarthNet2021}
formalized Earth-surface forecasting from Sentinel-2 imagery and weather
data. \href{https://arxiv.org/abs/2304.14065}{Presto} showed that a
small architecture specifically designed for multimodal time series can
compete with much larger models.

\textbf{Open scientific problems:}

\begin{itemize}
\tightlist
\item
  distinguishing semantic change from phenology, illumination,
  atmosphere, or sensor differences;
\item
  learning from irregular and incomplete time series;
\item
  handling abrupt rare events that are underrepresented in pretraining;
\item
  forecasting calibrated distributions rather than visually plausible
  single futures;
\item
  transferring temporal patterns across climate zones and agricultural
  calendars.
\end{itemize}

\hypertarget{generalization-domain-adaptation-and-limited-supervision}{%
\subsection{Generalization, domain adaptation, and limited
supervision}\label{generalization-domain-adaptation-and-limited-supervision}}

\textbf{Central objective:} Make models work where labels were not
collected.

EO distribution shifts include:

\begin{itemize}
\tightlist
\item
  \textbf{geographic shift:} another city, biome, country, or continent;
\item
  \textbf{temporal shift:} another season, year, climate regime, or
  post-disaster stage;
\item
  \textbf{sensor shift:} different spectral response functions,
  polarizations, noise, or product levels;
\item
  \textbf{resolution shift:} different ground-sampling distance or point
  density;
\item
  \textbf{acquisition shift:} view angle, illumination, orbit,
  atmosphere, or preprocessing;
\item
  \textbf{semantic shift:} different label ontology or class prevalence.
\end{itemize}

Current methods include transfer learning, unsupervised and source-free
domain adaptation, domain generalization, self-training, pseudo-labels,
few-shot learning, active learning, weak supervision, continual
learning, test-time adaptation, and parameter-efficient fine-tuning.

This branch is arguably the most important for real-world AI4EO. A
random chip split can place neighbouring, highly correlated samples into
train and test sets and thereby measure interpolation rather than
deployment transfer. Spatial and temporal separation must be part of the
scientific question.

Recent evidence is sobering:

\begin{itemize}
\tightlist
\item
  \href{https://arxiv.org/abs/2511.02831}{GeoCrossBench} reports large
  losses when models face satellites with non-overlapping bands or even
  additional unseen bands.
\item
  \href{https://arxiv.org/abs/2606.29664}{Shang et al.} found sharp
  regional degradation for foundation models in crop mapping.
\item
  \href{https://proceedings.mlr.press/v139/koh21a/koh21a.pdf}{WILDS},
  which includes the FMoW satellite benchmark, helped establish
  evaluation under naturally occurring distribution shift.
\end{itemize}

\textbf{Open scientific problem:} Better average in-domain accuracy can
coexist with worse performance on rare classes or unseen regions.
Robustness must be measured by geography, class, event, season, and
uncertainty---not only by one global score.

\hypertarget{trustworthy-uncertainty-aware-explainable-and-physics-aware-ai}{%
\subsection{Trustworthy, uncertainty-aware, explainable, and
physics-aware
AI}\label{trustworthy-uncertainty-aware-explainable-and-physics-aware-ai}}

\textbf{Central objective:} Produce outputs that scientists and
operational users can audit, calibrate, and reconcile with the
measurement process and Earth-system knowledge.

This branch contains several connected but distinct topics:

\begin{itemize}
\tightlist
\item
  aleatoric and epistemic uncertainty;
\item
  probability calibration and selective prediction;
\item
  out-of-distribution and anomaly detection;
\item
  interpretable features and explanation stability;
\item
  robustness to clouds, noise, corruption, and adversarial artefacts;
\item
  physical constraints, conservation relationships, and differentiable
  forward models;
\item
  hybrid numerical--ML models;
\item
  causal inference and process discovery;
\item
  provenance, traceability, fairness, privacy, and responsible use.
\end{itemize}

EO uncertainty enters at every stage: sensor noise, atmospheric
correction, geolocation, resampling, imperfect reference labels, model
parameters, and distribution shift. A 2024 study notes that no broadly
adopted, reliable, easy-to-use uncertainty framework yet exists for ML
in EO (\href{https://www.nature.com/articles/s41598-024-65954-w}{Singh
et al., 2024}). A 2025 benchmark specifically evaluates whether
pixel-wise uncertainty identifies segmentation errors and corrupted
regions (\href{https://arxiv.org/abs/2510.19586}{Rey et al., 2025}).

Physics-aware AI is especially relevant for geophysical retrieval and
downscaling. The objective is not merely to add a penalty called
``physics loss,'' but to identify a valid physical or measurement
constraint. Examples include:

\begin{itemize}
\tightlist
\item
  aggregate consistency between a high-resolution reconstruction and the
  original coarse observation;
\item
  non-negativity, boundedness, or conservation constraints;
\item
  sensor forward models and spectral response functions;
\item
  radiative-transfer constraints;
\item
  dynamics imposed by differential equations or numerical simulators.
\end{itemize}

\textbf{Open scientific problem:} Post-hoc visual explanations can look
convincing without being faithful, and a low test RMSE does not
guarantee physical validity or calibrated uncertainty under a new region
or event.

\hypertarget{vision-language-and-agentic-geospatial-ai}{%
\subsection{Vision-language and agentic geospatial
AI}\label{vision-language-and-agentic-geospatial-ai}}

\textbf{Central objective:} Connect EO measurements, spatial concepts,
scientific language, and executable tools.

The branch has three levels:

\begin{enumerate}
\def\labelenumi{\arabic{enumi}.}
\tightlist
\item
  \textbf{Image--text representation learning:} retrieval, zero-shot
  classification, and open-vocabulary recognition.
  \href{https://arxiv.org/abs/2306.11029}{RemoteCLIP} is a
  representative system.
\item
  \textbf{Conversational vision-language models:} captioning, VQA,
  grounding, counting, and multi-image reasoning. Representative systems
  include
  \href{https://openaccess.thecvf.com/content/CVPR2024/html/Kuckreja_GeoChat_Grounded_Large_Vision-Language_Model_for_Remote_Sensing_CVPR_2024_paper.html}{GeoChat}
  and \href{https://arxiv.org/abs/2412.15190}{EarthDial}.
\item
  \textbf{Tool-using EO agents:} select datasets and sensors, generate
  code, call catalogues or Earth Engine, execute analysis, inspect
  failures, and return evidence.
\end{enumerate}

This branch is scientifically interesting because EO analysis depends
heavily on metadata and procedural knowledge: correct collections,
units, scale factors, bands, cloud masks, projections, date ranges, and
spatial reducers.

It is also immature. The 2026
\href{https://aclanthology.org/2026.findings-acl.124.pdf}{UnivEARTH}
benchmark contains 408 evidence-grounded EO questions. In zero-shot
Google Earth Engine code generation, the best evaluated agent answered
only 40\% correctly and failed to execute code in more than 44\% of
cases. Reflection reduced syntax failures and raised accuracy toward
60\%, but correct execution still did not guarantee correct scientific
logic.

\textbf{Open scientific problems:}

\begin{itemize}
\tightlist
\item
  hallucinated collections, bands, units, dates, and scale factors;
\item
  failure to distinguish sensor products with similar names;
\item
  weak spatial and temporal reasoning;
\item
  answers that are plausible but unsupported by pixels or executable
  evidence;
\item
  poor uncertainty and provenance;
\item
  lack of formal verification for generated geospatial workflows.
\end{itemize}

\begin{center}\rule{0.5\linewidth}{0.5pt}\end{center}

\hypertarget{cross-cutting-enabling-layers}{%
\section{Cross-cutting enabling
layers}\label{cross-cutting-enabling-layers}}

\hypertarget{efficient-scalable-operational-and-onboard-ai}{%
\subsection{Efficient, scalable, operational, and onboard
AI}\label{efficient-scalable-operational-and-onboard-ai}}

EO systems operate at scales ranging from one local patch to the entire
planet and from ground-based clusters to satellite processors. Research
therefore includes:

\begin{itemize}
\tightlist
\item
  parameter-efficient fine-tuning;
\item
  knowledge distillation, pruning, quantization, and low-rank
  adaptation;
\item
  efficient data loading, tiling, and distributed inference;
\item
  learned compression and intelligent data selection;
\item
  low-latency disaster products;
\item
  edge and onboard inference;
\item
  energy, memory, bandwidth, and radiation constraints.
\end{itemize}

ESA's
\href{https://www.esa.int/Applications/Observing_the_Earth/Phsat-2}{\(\Phi\)sat-2},
launched on 16 August 2024, carries six AI applications including cloud
filtering, vessel detection/classification, image-to-map conversion,
compression, anomaly detection, and wildfire detection. A 2025 paper
reported the first on-orbit demonstration of a compressed geospatial
foundation model and emphasized that compression and domain adaptation
were both necessary for flight-ready inference
(\href{https://arxiv.org/abs/2512.01181}{Du et al.}).

This is a strategically important ESA-facing branch, but an MSc thesis
needs either access to representative hardware or a carefully justified
hardware simulator. FLOPs alone are not an operational evaluation.

\hypertarget{datasets-benchmarks-embeddings-and-representation-infrastructure}{%
\subsection{Datasets, benchmarks, embeddings, and representation
infrastructure}\label{datasets-benchmarks-embeddings-and-representation-infrastructure}}

This layer includes:

\begin{itemize}
\tightlist
\item
  AI-ready data cubes and harmonized products;
\item
  global pretraining corpora;
\item
  label quality and ontology alignment;
\item
  spatially and temporally explicit splits;
\item
  benchmark harnesses and reproduced baselines;
\item
  model cards, data cards, licences, and provenance;
\item
  global embedding layers and similarity search;
\item
  STAC catalogues, cloud-native formats, and reproducible pipelines.
\end{itemize}

The 2026 audit by \href{https://arxiv.org/html/2605.12678v1}{Corley et
al.} quantified the problem: among 152 audited papers, it found 46
cross-paper disagreements of at least 10 points for nominally the same
model, benchmark, and protocol; 94 of 126 papers with extractable
pretraining configurations used a configuration no other paper used; and
39\% released no weights.

Benchmark and data work is therefore not secondary administration. It is
a substantive scientific branch when it establishes a capability that
was previously impossible to measure.

\begin{center}\rule{0.5\linewidth}{0.5pt}\end{center}

\hypertarget{the-principal-unresolved-ai4eo-problems}{%
\section{The principal unresolved AI4EO
problems}\label{the-principal-unresolved-ai4eo-problems}}

\begin{longtable}[]{@{}
  >{\raggedright\arraybackslash}p{(\columnwidth - 4\tabcolsep) * \real{0.3333}}
  >{\raggedright\arraybackslash}p{(\columnwidth - 4\tabcolsep) * \real{0.3333}}
  >{\raggedright\arraybackslash}p{(\columnwidth - 4\tabcolsep) * \real{0.3333}}@{}}
\toprule\noalign{}
\begin{minipage}[b]{\linewidth}\raggedright
Problem
\end{minipage} & \begin{minipage}[b]{\linewidth}\raggedright
Why it is specifically difficult in EO
\end{minipage} & \begin{minipage}[b]{\linewidth}\raggedright
Researchable question
\end{minipage} \\
\midrule\noalign{}
\endhead
\bottomrule\noalign{}
\endlastfoot
\textbf{Geographic generalization} & Nearby pixels are correlated;
land-cover appearance and class prevalence vary by region & Can
adaptation improve leave-one-region-out performance without harming rare
classes? \\
\textbf{Temporal and climate shift} & Phenology, land management,
atmosphere, and extremes change across seasons and years & Which
temporal representations remain stable across years and climate
zones? \\
\textbf{Cross-sensor transfer} & Sensors have different spectral
response functions, polarizations, resolutions, and noise & Can
wavelength/metadata-aware adapters transfer between instruments with
partially overlapping bands? \\
\textbf{Label scarcity and noise} & Expert labels are expensive,
temporally stale, spatially imprecise, and ontology-dependent & Which
SSL, active-learning, or weak-supervision method gives the best gain per
labelled area? \\
\textbf{Rare events and long tails} & Extreme fires, floods, damage
levels, and minority crops are scientifically important but
statistically rare & Can event-aware sampling and calibrated uncertainty
prevent majority-class collapse? \\
\textbf{Multimodal alignment} & Acquisitions are asynchronous,
co-registration is imperfect, and missingness is not random & Does
fusion remain useful under realistic cloud, delay, and sensor-outage
patterns? \\
\textbf{Resolution and scale} & Resampling changes grid size but not
native information content; objects appear at radically different scales
& Can models represent native resolutions without false detail or
excessive compute? \\
\textbf{Benchmark leakage} & Random chip splits can share scenes,
neighbourhoods, dates, or derived labels across partitions & How much
reported performance disappears under spatially and temporally blocked
evaluation? \\
\textbf{Foundation-model comparability} & Input bands, decoders,
preprocessing, adaptation, and metrics vary across papers & Which
capabilities survive a shared data and adaptation protocol? \\
\textbf{Uncertainty and OOD behaviour} & Operational decisions need to
know where a map is unreliable & Which uncertainty method remains
calibrated under geographic and sensor shift? \\
\textbf{Physical faithfulness} & A visually plausible output may violate
the sensor observation or Earth-system constraints & Can explicit
forward/aggregation consistency prevent hallucinated fine-scale
structure? \\
\textbf{Generative reliability} & Cross-modal generation has many
plausible outputs, but only one Earth was observed & How should
conditional distributions, ambiguity, and downstream bias be
evaluated? \\
\textbf{Agent reliability} & EO tools require exact product IDs, bands,
scale factors, projections, and temporal logic & Can typed tools,
retrieval, and automatic scientific checks make generated workflows
auditable? \\
\textbf{Operational efficiency} & Planet-scale inference and onboard
processing have strict storage, energy, and latency limits & What is the
accuracy--calibration--latency frontier after compression? \\
\textbf{Data and representation inequality} & Training data and
benchmarks overrepresent some continents, climates, and commercial data
users & How does performance vary by geography, and which sampling
strategy closes the gap? \\
\end{longtable}

\hypertarget{why-spatially-correct-evaluation-is-central}{%
\subsection{Why spatially correct evaluation is
central}\label{why-spatially-correct-evaluation-is-central}}

Suppose adjacent patches from one Sentinel-2 tile are randomly divided
into training and test sets. They share acquisition conditions, local
land-cover structure, atmospheric state, preprocessing, and often label
sources. A high test score may therefore reflect local interpolation.

A deployment-oriented protocol should separate data according to the
intended claim:

\begin{itemize}
\tightlist
\item
  \textbf{new place:} spatial blocks, countries, cities, watersheds, or
  events;
\item
  \textbf{new time:} later year or unseen season;
\item
  \textbf{new sensor:} held-out platform or band configuration;
\item
  \textbf{few labels:} explicit n-shot or labelled-area budget;
\item
  \textbf{rare event:} leave-one-event-out or leave-one-hazard-out;
\item
  \textbf{operational uncertainty:} calibration and risk--coverage under
  the same shift.
\end{itemize}

The split is part of the hypothesis, not a bookkeeping decision.

\hypertarget{why-best-foundation-model-is-currently-the-wrong-question}{%
\subsection{Why ``best foundation model'' is currently the wrong
question}\label{why-best-foundation-model-is-currently-the-wrong-question}}

The more useful questions are:

\begin{itemize}
\tightlist
\item
  best for \textbf{which sensor and band set};
\item
  best under \textbf{which label budget};
\item
  best with \textbf{frozen, parameter-efficient, or full fine-tuning};
\item
  best for \textbf{in-domain accuracy or cross-region transfer};
\item
  best at \textbf{which spatial resolution and map scale};
\item
  best with \textbf{which uncertainty and compute requirements}.
\end{itemize}

Model claims should be capability-specific. TerraMind, OlmoEarth, and
RAMEN report very strong results, but their own papers use different
data, adaptation choices, and benchmark subsets. Their progress is real;
a universal total ordering is not yet supported.

\begin{center}\rule{0.5\linewidth}{0.5pt}\end{center}

\hypertarget{current-scientific-and-agency-interests}{%
\section{Current scientific and agency
interests}\label{current-scientific-and-agency-interests}}

\hypertarget{esa-facing-interests}{%
\subsection{ESA-facing interests}\label{esa-facing-interests}}

Recent ESA \(\Phi\)-lab work and calls indicate sustained interest in:

\begin{itemize}
\tightlist
\item
  multimodal EO foundation models and benchmarking;
\item
  trustworthy, explainable, and physics-aware AI;
\item
  compact foundation models and onboard intelligence;
\item
  disaster response and extreme-event generalization;
\item
  global embeddings and AI-ready data infrastructure;
\item
  natural-language and code-based access to EO;
\item
  digital twins and integration with Earth-system models.
\end{itemize}

TerraMind is a major current ESA/IBM programme, while PANGAEA and PhilEO
illustrate the emphasis on independent evaluation. \(\Phi\)sat-2 makes
compression, raw-to-product processing, application updating, and
onboard inference operational research rather than speculative topics.
ESA's 2025 Living Planet Symposium explicitly highlighted explainable,
trustworthy, and physics-aware AI
(\href{https://philab.esa.int/unmissable-\%CF\%86-lab-moments-at-lps-2025/}{ESA
\(\Phi\)-lab}).

\hypertarget{nasa-facing-interests}{%
\subsection{NASA-facing interests}\label{nasa-facing-interests}}

NASA's open-science foundation-model programme includes:

\begin{itemize}
\tightlist
\item
  Prithvi-EO for HLS imagery;
\item
  Prithvi Weather--Climate for MERRA-2 atmospheric data;
\item
  open models, code, and downstream scientific applications;
\item
  disaster response, crop and land-use mapping, ecosystem monitoring,
  and LST;
\item
  model deployment and collaboration with domain scientists.
\end{itemize}

\href{https://arxiv.org/abs/2412.02732}{Prithvi-EO-2.0} was trained on
4.2 million global HLS time-series samples and released in 300M and 600M
variants. NASA describes the model as part of a wider open family of
scientific foundation models
(\href{https://science.nasa.gov/science-research/ai-geospatial-model-earth/}{NASA,
2024}).

\hypertarget{application-domains-with-sustained-research-demand}{%
\subsection{Application domains with sustained research
demand}\label{application-domains-with-sustained-research-demand}}

\begin{itemize}
\tightlist
\item
  \textbf{Disasters:} floods, wildfire, earthquake damage, landslides,
  rapid mapping, and recovery.
\item
  \textbf{Agriculture and food security:} crop type, field boundaries,
  phenology, yield, irrigation, and transfer to data-scarce regions.
\item
  \textbf{Forestry and biodiversity:} biomass, canopy height,
  deforestation, habitat, species distributions, and ecosystem change.
\item
  \textbf{Climate and urban systems:} LST, urban heat, emissions,
  downscaling, extremes, and adaptation.
\item
  \textbf{Water, ocean, and cryosphere:} surface water, soil moisture,
  coastal change, sea ice, water quality, and marine anomalies.
\item
  \textbf{Atmosphere and weather:} retrieval, nowcasting, data
  assimilation, parameterization, and hybrid numerical--AI models.
\end{itemize}

Applications alone do not define novelty. ``Flood mapping with a
Transformer'' is likely too broad and incremental. ``Calibrated
leave-one-event-out flood mapping with missing optical observations''
defines a research contribution.

\begin{center}\rule{0.5\linewidth}{0.5pt}\end{center}

\hypertarget{designing-a-defensible-msc-thesis}{%
\section{Designing a defensible MSc
thesis}\label{designing-a-defensible-msc-thesis}}

\hypertarget{thesis-formula}{%
\subsection{Thesis formula}\label{thesis-formula}}

A strong AI4EO MSc thesis can usually be expressed as:

\begin{quote}
\textbf{For {[}EO phenomenon/task{]}, under {[}realistic shift or data
constraint{]}, determine whether {[}methodological contribution{]}
improves {[}scientifically appropriate metrics{]} relative to {[}strong
classical, supervised, and pretrained baselines{]}.}
\end{quote}

Every candidate should have:

\begin{enumerate}
\def\labelenumi{\arabic{enumi}.}
\tightlist
\item
  \textbf{A precise failure mode}, not merely a new architecture.
\item
  \textbf{Open or institutionally guaranteed data.}
\item
  \textbf{A leakage-resistant split defined before modelling.}
\item
  \textbf{At least three baseline families:} classical/physical,
  task-specific deep model, and pretrained model.
\item
  \textbf{Metrics matched to the claim:} accuracy alone is rarely
  sufficient.
\item
  \textbf{A compute plan that uses pretrained weights rather than new
  foundation-model pretraining.}
\item
  \textbf{A negative-result path:} the thesis remains valuable if the
  foundation model does not win.
\end{enumerate}

\hypertarget{feasibility-gates}{%
\subsection{Feasibility gates}\label{feasibility-gates}}

Reject or redesign a candidate if any of the following is unresolved:

\begin{itemize}
\tightlist
\item
  The reference labels cannot support the claimed spatial resolution.
\item
  Train and test data cannot be separated by the intended
  geography/time/event.
\item
  The proposed novelty is only replacing a CNN with a Transformer.
\item
  Required imagery is commercial or access is uncertain.
\item
  The experiment requires pretraining at institutional-cluster scale.
\item
  The only evaluation is one saturated dataset with random patches.
\item
  The contribution cannot be isolated through ablations.
\item
  ``Physical consistency'' is asserted without an explicit measurement
  or process constraint.
\end{itemize}

\begin{center}\rule{0.5\linewidth}{0.5pt}\end{center}

\begin{landscape}
\hypertarget{candidate-thesis-directions-tailored-to-stefanos-background}{%
\section{Candidate thesis directions tailored to Stefano's
background}\label{candidate-thesis-directions-tailored-to-stefanos-background}}

\hypertarget{comparative-decision-matrix}{%
\subsection{Comparative decision
matrix}\label{comparative-decision-matrix}}

\begin{footnotesize}
\setlength{\tabcolsep}{3pt}
\begin{longtable}[]{@{}
  >{\raggedright\arraybackslash}p{(\linewidth - 14\tabcolsep) * \real{0.2300}}
  >{\raggedright\arraybackslash}p{(\linewidth - 14\tabcolsep) * \real{0.2000}}
  >{\centering\arraybackslash}p{(\linewidth - 14\tabcolsep) * \real{0.0700}}
  >{\centering\arraybackslash}p{(\linewidth - 14\tabcolsep) * \real{0.0700}}
  >{\centering\arraybackslash}p{(\linewidth - 14\tabcolsep) * \real{0.1000}}
  >{\centering\arraybackslash}p{(\linewidth - 14\tabcolsep) * \real{0.0900}}
  >{\raggedright\arraybackslash}p{(\linewidth - 14\tabcolsep) * \real{0.1500}}
  >{\centering\arraybackslash}p{(\linewidth - 14\tabcolsep) * \real{0.0900}}@{}}
\toprule\noalign{}
\begin{minipage}[b]{\linewidth}\raggedright
Candidate
\end{minipage} & \begin{minipage}[b]{\linewidth}\raggedright
Core contribution
\end{minipage} & \begin{minipage}[b]{\linewidth}\raggedleft
Open data
\end{minipage} & \begin{minipage}[b]{\linewidth}\raggedleft
Compute
\end{minipage} & \begin{minipage}[b]{\linewidth}\raggedleft
Evaluation clarity
\end{minipage} & \begin{minipage}[b]{\linewidth}\raggedleft
Novelty potential
\end{minipage} & \begin{minipage}[b]{\linewidth}\raggedright
Principal risk
\end{minipage} & \begin{minipage}[b]{\linewidth}\raggedleft
Provisional fit
\end{minipage} \\
\midrule\noalign{}
\endhead
\bottomrule\noalign{}
\endlastfoot
\textbf{A. Uncertainty-aware, physically consistent LST downscaling with
EO-FM features} & Cross-region thermal sharpening with aggregation
consistency and calibrated uncertainty & High & Medium & High if
validation is designed carefully & High & No true 10 m thermal ground
truth & \textbf{Very high} \\
\textbf{B. Geographic transfer and parameter-efficient adaptation of EO
foundation models} & Controlled comparison under spatial/temporal
shifts; adapters or domain-distance-aware selection & High & Medium &
Very high & Medium--high & Can become ``just a benchmark'' without a
method contribution & \textbf{Very high} \\
\textbf{C. Missing-modality robust Sentinel-1/2 fusion} & Fusion that
degrades gracefully under clouds, delay, and sensor absence & High &
Medium & High & High & Simulated missingness may be unrealistic &
\textbf{High} \\
\textbf{D. Uncertainty-aware disaster change detection under event
shift} & Leave-one-event-out change maps with calibrated pixel
uncertainty & High & Medium & High & High & Labels across events are
heterogeneous/noisy & \textbf{High} \\
\textbf{E. Sensor-aware, verifiable EO agent} & Retrieval + typed
constraints + execution checks for auditable EO workflows & Medium--high
& Low--medium & High with UnivEARTH-style benchmark & Very high &
Rapidly moving field; API/tool dependence & \textbf{High but riskier} \\
\textbf{F. Compact/onboard adaptation of an EO foundation model} &
Quantization/distillation and accuracy--latency--calibration frontier &
High & Medium & Medium--high with hardware & High & Weak without
representative hardware & \textbf{Conditional} \\
\textbf{G. Cross-sensor hyperspectral representation transfer} &
Spectral-response-aware adaptation across hyperspectral/multispectral
sensors & Medium & Medium & Medium & High & Labels and realistic
cross-sensor pairs are scarce & \textbf{Conditional} \\
\end{longtable}
\end{footnotesize}
\end{landscape}

\hypertarget{candidate-a-lst-downscaling-with-uncertainty-and-physical-consistency}{%
\subsection{Candidate A --- LST downscaling with uncertainty and
physical
consistency}\label{candidate-a-lst-downscaling-with-uncertainty-and-physical-consistency}}

\textbf{Working title}\\
\emph{Uncertainty-Aware and Physically Consistent
Land-Surface-Temperature Downscaling Using Earth-Observation
Foundation-Model Representations}

\textbf{Why it fits}

\begin{itemize}
\tightlist
\item
  It builds on prior Landsat 8/9 and Sentinel-2 LST work rather than
  discarding that investment.
\item
  It adds current AI4EO questions: pretrained representations,
  geographic transfer, physical consistency, and uncertainty.
\item
  It connects image/signal processing, regression, multimodal fusion,
  and environmental application.
\item
  It can produce a useful negative result if pretrained features do not
  outperform carefully tuned RF/CNN baselines.
\end{itemize}

\textbf{Possible research question}

\begin{quote}
Do frozen or parameter-efficient EO foundation-model features improve
cross-city and cross-season LST downscaling, and can coarse-scale
consistency plus uncertainty calibration prevent unsupported fine-scale
thermal detail?
\end{quote}

\textbf{Key methodological contribution}

Predict high-resolution residual structure using Sentinel-2, land-cover,
DEM, and possibly weather/context features, while enforcing that the
downscaled prediction aggregates back to the observed coarse thermal
measurement:

\[
\mathcal{L} = \mathcal{L}_{\text{reconstruction}} + \lambda_{c}\,\left\|D(\hat{T}_{HR})-T_{LR}\right\|_1 + \lambda_{u}\,\mathcal{L}_{\text{uncertainty}},
\]

where \(D\) is the sensor-aware degradation or aggregation operator,
\(T_{LR}\) is observed coarse LST, and \(\hat{T}_{HR}\) is the
predicted high-resolution field.

\textbf{Essential caution}

Downscaling cannot manufacture measured 10 m thermal information.
Because true 10 m LST ground truth is normally unavailable, a credible
thesis must combine several validation levels:

\begin{enumerate}
\def\labelenumi{\arabic{enumi}.}
\tightlist
\item
  reduced-resolution simulation where the target is observable;
\item
  coarse-scale reconstruction/energy consistency;
\item
  cross-city and cross-season transfer;
\item
  independent comparison with temporally matched thermal products or
  stations where available;
\item
  spatial uncertainty and risk--coverage analysis;
\item
  tests of whether added fine detail tracks independent urban morphology
  rather than merely sharpening optical edges.
\end{enumerate}

\textbf{Baselines}

\begin{itemize}
\tightlist
\item
  linear regression and Random Forest;
\item
  a task-specific CNN/U-Net or geographically weighted model;
\item
  frozen features from one or two open EO encoders such as Prithvi,
  DOFA, or a TerraMind compact variant;
\item
  ablations without consistency and without uncertainty.
\end{itemize}

\textbf{Novelty boundary}

IBM already provides a
\href{https://github.com/ibm-granite/granite-geospatial-land-surface-temperature}{foundation-model-based
LST system}, and Landsat--Sentinel thermal fusion remains active
research. The thesis novelty should therefore be \textbf{cross-domain
reliability and physical validation}, not simply ``apply Prithvi to
LST.''

\hypertarget{candidate-b-geographic-transfer-and-efficient-adaptation}{%
\subsection{Candidate B --- Geographic transfer and efficient
adaptation}\label{candidate-b-geographic-transfer-and-efficient-adaptation}}

\textbf{Working title}\\
\emph{Capability-Oriented Evaluation and Parameter-Efficient Adaptation
of Earth-Observation Foundation Models under Geographic Shift}

\textbf{Research question}

\begin{quote}
Under spatially disjoint evaluation, which adaptation strategies
preserve rare-class performance and calibration, and can a measurable
domain-distance signal predict which model or adapter will transfer
best?
\end{quote}

\textbf{Contribution options}

\begin{itemize}
\tightlist
\item
  compare frozen, linear-probe, adapter/LoRA, and full fine-tuning under
  equal compute;
\item
  quantify target-domain distance using embeddings, metadata, or optimal
  transport;
\item
  select the model/adapter before target labels are fully available;
\item
  report class-wise transfer, calibration, and compute rather than only
  average accuracy.
\end{itemize}

\textbf{Data and tools}

\begin{itemize}
\tightlist
\item
  one PANGAEA or GEO-Bench-2 task with genuine geographic units;
\item
  BigEarthNet v2 country/region splits;
\item
  Sen1Floods11 leave-one-event-out;
\item
  an agricultural benchmark with region-separated labels.
\end{itemize}

\textbf{Baselines}

\begin{itemize}
\tightlist
\item
  Random Forest or gradient boosting on spectral/temporal features;
\item
  ResNet/U-Net or a small ViT trained from scratch;
\item
  two or three open GFMs selected for complementary capabilities, not a
  long leaderboard.
\end{itemize}

\textbf{What makes it a thesis rather than a benchmark report}

Add a method or scientific hypothesis---for example,
domain-distance-guided adapter selection, rare-class-aware adaptation,
or calibration-aware early stopping.

\hypertarget{candidate-c-robust-opticalsar-fusion-with-missing-observations}{%
\subsection{Candidate C --- Robust optical--SAR fusion with missing
observations}\label{candidate-c-robust-opticalsar-fusion-with-missing-observations}}

\textbf{Working title}\\
\emph{Robust Multimodal Earth Observation under Cloud, Delay, and Sensor
Absence}

\textbf{Research question}

\begin{quote}
Can modality-dropout, asynchronous temporal encoding, or
uncertainty-aware fusion retain useful performance when optical and SAR
observations are not perfectly paired?
\end{quote}

\textbf{Suitable data}

\begin{itemize}
\tightlist
\item
  \href{https://arxiv.org/abs/2201.09613}{SEN12MS-CR-TS} for cloudy
  optical, clear optical, and SAR time series;
\item
  BigEarthNet-MM for paired Sentinel-1/2 classification;
\item
  PASTIS-R for agricultural radar--optical time series.
\end{itemize}

\textbf{Evaluation regimes}

\begin{itemize}
\tightlist
\item
  all modalities present;
\item
  optical missing at random;
\item
  optical missing conditional on real cloud masks;
\item
  asynchronous observation delays;
\item
  sensor outage at test time;
\item
  transfer to an unseen geography or season.
\end{itemize}

\textbf{Metrics}

Task performance versus missingness rate, calibration, worst-group/event
performance, inference cost, and degradation relative to the
full-modality upper bound.

\textbf{Main danger}

Artificially deleting channels uniformly is not a realistic missing-data
model. Missingness should follow cloud, revisit, or acquisition
patterns.

\hypertarget{candidate-d-uncertainty-aware-change-or-disaster-mapping}{%
\subsection{Candidate D --- Uncertainty-aware change or disaster
mapping}\label{candidate-d-uncertainty-aware-change-or-disaster-mapping}}

\textbf{Working title}\\
\emph{Calibrated Semantic Change Detection under Unseen Disaster Events}

\textbf{Research question}

\begin{quote}
Which uncertainty method best identifies incorrect change predictions
when the test event, geography, or sensor conditions differ from
training?
\end{quote}

\textbf{Experimental design}

\begin{itemize}
\tightlist
\item
  choose one hazard and one clear output, such as flood extent or
  burned-area segmentation;
\item
  train on several events and hold out complete events;
\item
  compare deterministic confidence, deep ensembles,
  evidential/probabilistic heads, and conformal risk control;
\item
  evaluate IoU/F1 alongside Brier score, expected calibration error,
  negative log-likelihood, risk--coverage, and error-detection AUROC;
\item
  analyse boundaries, rare severity levels, and corrupted/missing inputs
  separately.
\end{itemize}

\textbf{Why it is useful}

Emergency mapping needs a map of \emph{where the model may be wrong},
not only a globally averaged accuracy score.

\hypertarget{candidate-e-verifiable-eo-agents}{%
\subsection{Candidate E --- Verifiable EO
agents}\label{candidate-e-verifiable-eo-agents}}

\textbf{Working title}\\
\emph{Sensor-Aware Verification for Tool-Using Earth-Observation Agents}

\textbf{Research question}

\begin{quote}
Can retrieval of authoritative metadata, typed tool schemas, and
automatic physical/unit checks reduce execution and scientific-logic
failures in EO code-generating agents?
\end{quote}

\textbf{Possible contribution}

\begin{itemize}
\tightlist
\item
  retrieve product-specific documentation before code generation;
\item
  constrain collection IDs, band names, temporal availability, units,
  and scale factors;
\item
  run static checks and small sentinel queries before full execution;
\item
  verify numerical ranges and provenance;
\item
  compare zero-shot, retrieval-only, reflection-only, and
  constrained-agent conditions on UnivEARTH.
\end{itemize}

\textbf{Why it fits}

It uses Stefano's NLP and multi-agent background while remaining
grounded in actual EO data and sensor knowledge.

\textbf{Why it is riskier}

The field changes quickly; cloud APIs and base models may change during
the thesis; reproducibility needs cached questions, pinned tools, and
recorded outputs.

\hypertarget{candidate-f-efficient-and-onboard-eo-foundation-models}{%
\subsection{Candidate F --- Efficient and onboard EO foundation
models}\label{candidate-f-efficient-and-onboard-eo-foundation-models}}

\textbf{Research question}

\begin{quote}
What accuracy--calibration--latency trade-off is achieved by
distillation or quantization of an EO encoder under realistic onboard
inputs and distribution shift?
\end{quote}

This becomes strong only if representative hardware, an emulator, or
collaboration is available. Otherwise it risks becoming a generic
compression comparison with estimated FLOPs. ESA \(\Phi\)sat-2 makes it
strategically relevant, but operational constraints must be real.

\hypertarget{candidate-g-cross-sensor-hyperspectral-transfer}{%
\subsection{Candidate G --- Cross-sensor hyperspectral
transfer}\label{candidate-g-cross-sensor-hyperspectral-transfer}}

\textbf{Research question}

\begin{quote}
Can wavelength-conditioned representations transfer material or crop
classifiers across sensors with different band centres and spectral
response functions?
\end{quote}

This is scientifically attractive and links to recent hyperspectral
study. A credible version should avoid the standard random-pixel Indian
Pines experiment. It needs spatially disjoint regions, preferably more
than one sensor, explicit spectral-response modelling, and limited-label
evaluation. Data availability and ground truth should be confirmed
before selecting it.

\begin{center}\rule{0.5\linewidth}{0.5pt}\end{center}

\hypertarget{provisional-recommendation}{%
\section{Provisional recommendation}\label{provisional-recommendation}}

The directions can be grouped by thesis character:

\begin{longtable}[]{@{}
  >{\raggedright\arraybackslash}p{(\columnwidth - 2\tabcolsep) * \real{0.5000}}
  >{\raggedright\arraybackslash}p{(\columnwidth - 2\tabcolsep) * \real{0.5000}}@{}}
\toprule\noalign{}
\begin{minipage}[b]{\linewidth}\raggedright
Desired thesis character
\end{minipage} & \begin{minipage}[b]{\linewidth}\raggedright
Best candidate
\end{minipage} \\
\midrule\noalign{}
\endhead
\bottomrule\noalign{}
\endlastfoot
Strong continuity with earlier work and a clear environmental variable &
\textbf{A --- LST downscaling} \\
Most methodologically general ML/CV thesis & \textbf{B --- geographic
transfer/adaptation} \\
Strong multisensor EO identity & \textbf{C --- missing-modality S1/S2
fusion} \\
Strong operational/trustworthy-AI identity & \textbf{D ---
uncertainty-aware disaster mapping} \\
Strong NLP/agentic and emerging-topic identity & \textbf{E ---
verifiable EO agents} \\
Strong ESA upstream/edge identity, if hardware is available & \textbf{F
--- onboard adaptation} \\
Strong imaging-spectroscopy identity, if suitable paired data exist &
\textbf{G --- hyperspectral transfer} \\
\end{longtable}

The recommended discussion order with the professors is:

\begin{enumerate}
\def\labelenumi{\arabic{enumi}.}
\tightlist
\item
  \textbf{Candidate A}, because it combines existing domain experience
  with several open AI4EO problems and can be scoped from conservative
  to ambitious.
\item
  \textbf{Candidate B}, because it has the cleanest methodology and the
  lowest dependence on a particular Earth-science application.
\item
  \textbf{Candidate C or D}, depending on whether the group prefers
  multisensor representation learning or trustworthy disaster mapping.
\item
  \textbf{Candidate E}, as a higher-risk, high-novelty alternative that
  uses the NLP background.
\end{enumerate}

The immediate decision should not be ``which architecture?'' It should
be:

\begin{itemize}
\tightlist
\item
  Which scientific variable or operational task matters to the
  supervisors?
\item
  Which geographic/temporal transfer claim can be evaluated?
\item
  Which labels and compute are guaranteed?
\item
  Does the group prefer a method contribution, an application
  contribution, or a benchmark/reliability contribution?
\end{itemize}

\begin{center}\rule{0.5\linewidth}{0.5pt}\end{center}

\hypertarget{recommended-reading-sequence}{%
\section{Recommended reading
sequence}\label{recommended-reading-sequence}}

\hypertarget{stage-1-understand-the-fields-development}{%
\subsection{Stage 1 --- Understand the field's
development}\label{stage-1-understand-the-fields-development}}

\begin{enumerate}
\def\labelenumi{\arabic{enumi}.}
\item
  \href{https://arxiv.org/abs/1710.03959}{Zhu et al.~--- Deep Learning
  in Remote Sensing: A Comprehensive Review and List of Resources
  (2017)}\\
  Read for the transition from classical remote sensing to learned
  spectral--spatial representations.
\item
  \href{https://arxiv.org/html/2305.08413v2}{Tuia et al.~--- Artificial
  Intelligence to Advance Earth Observation: A Perspective
  (2023/2024)}\\
  Read for the wider progression from mapping to explanation, physics,
  communication, trust, and ethics.
\item
  \href{https://arxiv.org/html/2312.05327v3}{Schmitt et al.~---
  Data-centric Machine Learning for Geospatial Remote Sensing Data
  (2024)}\\
  Read for label quality, geographic diversity, data curation, and
  evaluation.
\end{enumerate}

\hypertarget{stage-2-understand-representation-learning}{%
\subsection{Stage 2 --- Understand representation
learning}\label{stage-2-understand-representation-learning}}

\begin{enumerate}
\def\labelenumi{\arabic{enumi}.}
\setcounter{enumi}{3}
\item
  \href{https://openaccess.thecvf.com/content/ICCV2021/papers/Manas_Seasonal_Contrast_Unsupervised_Pre-Training_From_Uncurated_Remote_Sensing_Data_ICCV_2021_paper.pdf}{Mañas
  et al.~--- Seasonal Contrast / SeCo (ICCV 2021)}\\
  Why temporal consistency is an EO-specific self-supervision signal.
\item
  \href{https://proceedings.neurips.cc/paper_files/paper/2022/hash/01c561df365429f33fcd7a7faa44c985-Abstract-Conference.html}{Cong
  et al.~--- SatMAE (NeurIPS 2022)}\\
  How masked modelling is adapted to spectral and temporal structure.
\item
  \href{https://arxiv.org/abs/2311.00566}{Fuller et al.~--- CROMA
  (2023)}\\
  Radar--optical contrastive and masked pretraining.
\item
  \href{https://arxiv.org/abs/2403.15356}{Xiong et al.~--- DOFA
  (2024/2025)}\\
  Cross-sensor modelling with wavelength-conditioned parameters.
\item
  \href{https://arxiv.org/abs/2412.14123}{Astruc et al.~--- AnySat (CVPR
  2025)}\\
  Heterogeneous resolutions, scales, modalities, and datasets in one
  model.
\item
  \href{https://arxiv.org/html/2504.11171v5}{Jakubik et al.~---
  TerraMind (2025/2026)}\\
  Generative multimodality and the current ESA/IBM direction.
\end{enumerate}

\hypertarget{stage-3-understand-what-is-still-unsolved}{%
\subsection{Stage 3 --- Understand what is still
unsolved}\label{stage-3-understand-what-is-still-unsolved}}

\begin{enumerate}
\def\labelenumi{\arabic{enumi}.}
\setcounter{enumi}{9}
\item
  \href{https://proceedings.neurips.cc/paper_files/paper/2023/hash/a0644215d9cff6646fa334dfa5d29c5a-Abstract-Datasets_and_Benchmarks.html}{Lacoste
  et al.~--- GEO-Bench (NeurIPS 2023)}\\
  Standardized multi-task evaluation.
\item
  \href{https://arxiv.org/abs/2412.04204}{Marsocci et al.~--- PANGAEA
  (2024/2025)}\\
  Why GFMs do not universally beat supervised baselines.
\item
  \href{https://arxiv.org/html/2511.15658v1}{GEO-Bench-2 (2025)}\\
  Capability-oriented evaluation and the distinction between RGB and
  EO-specific advantages.
\item
  \href{https://arxiv.org/html/2605.12678v1}{Corley et al.~--- No One
  Knows the State of the Art in Geospatial Foundation Models (2026)}\\
  Essential reading before accepting any leaderboard or ``best model''
  claim.
\item
  \href{https://arxiv.org/abs/2606.29664}{Shang et al.~--- Benchmarking
  GFMs for Agriculture Applications (2026)}\\
  A concrete example of geographic transfer failure.
\item
  \href{https://aclanthology.org/2026.findings-acl.124.pdf}{Kao et
  al.~--- Towards LLM Agents for Earth Observation (ACL 2026)}\\
  Evidence that agentic EO is promising but far from reliable.
\end{enumerate}

For each paper, use the established review template:

\begin{itemize}
\tightlist
\item
  problem;
\item
  why it matters;
\item
  dataset and geographic coverage;
\item
  model and input modalities;
\item
  claimed contribution;
\item
  evaluation split and metrics;
\item
  results;
\item
  weaknesses and missing baselines;
\item
  relevance to a possible thesis.
\end{itemize}

\begin{center}\rule{0.5\linewidth}{0.5pt}\end{center}

\hypertarget{questions-to-take-to-the-thesis-meeting}{%
\section{Questions to take to the thesis
meeting}\label{questions-to-take-to-the-thesis-meeting}}

\begin{enumerate}
\def\labelenumi{\arabic{enumi}.}
\tightlist
\item
  Should the thesis make its main contribution in \textbf{AI
  methodology}, \textbf{EO application}, or
  \textbf{evaluation/reliability}?
\item
  Is the group interested in continuing the LST/urban-heat work, or
  should the scientific application change?
\item
  What GPU, storage, and cloud/HPC resources are guaranteed for the full
  thesis period?
\item
  Does the group have access to expert labels, in-situ measurements, or
  a domain collaborator?
\item
  Would a negative result about foundation-model transfer be considered
  a valid contribution if the evaluation is rigorous?
\item
  Is collaboration with an ESA/EO group realistic, or should the thesis
  be fully reproducible from open data and weights?
\item
  Which deployment claim matters: another city, another season, another
  sensor, another event, or low-label adaptation?
\end{enumerate}

\begin{center}\rule{0.5\linewidth}{0.5pt}\end{center}

\hypertarget{final-assessment}{%
\section{Final assessment}\label{final-assessment}}

AI4EO has moved beyond the simple transplantation of computer-vision
architectures to satellite images. Its scientific frontier is defined by
the structure of EO itself: many sensing physics, multiple native
resolutions, repeated but irregular observations, global spatial
heterogeneity, scarce and noisy labels, rare extremes, and the
requirement that outputs remain scientifically and operationally
credible.

Foundation models are an important part of this frontier, but they are
not the whole field and not automatically the best solution. The most
valuable current research asks what capabilities pretraining creates,
where those capabilities fail, how they can be adapted efficiently, and
how their uncertainty and physical consistency can be demonstrated
outside the benchmark on which they were tuned.

For an MSc thesis, the opportunity is therefore to choose one precise
intersection of:

\begin{enumerate}
\def\labelenumi{\arabic{enumi}.}
\tightlist
\item
  an Earth-observation problem;
\item
  an observable and reproducible dataset;
\item
  a real deployment shift or constraint;
\item
  a methodological contribution;
\item
  a rigorous evaluation protocol.
\end{enumerate}

That structure makes the work useful even if the newest model does not
win---which is exactly what a scientifically strong thesis should allow.

\end{document}
